% META: {"id":"1SPE-VARALEA-CO-028","chapitre":"1SPE-VARIABLES-ALEATOIRES","type_objet":"corrige","exercice_ref":"1SPE-VARALEA-EX-028","capacites_codes":[],"extension_codes":["X1"],"programme_alignment":"OPTIONAL_EXTENSION","extension_label":"Approfondissement — Vers la Terminale","sources_inspiration":[],"status":"approved","origin":"generated","status_history":["generated","audited","accepted","approved"],"release_acceptance":"RELEASE_OWNER_BATCH_ACCEPTANCE","acceptance_closure_digest":"sha256:9b3ccf9a81c5520fbb7e03b7b2d3e7bf2057a02834b6d908bf3be5f49f3b553f","acceptance_receipt_digest":"sha256:4fad86f0282e0fa4e0419cca1ad6379ae7af5a280c04a4a90880f90a1c044242"}
\begin{center}\textbf{Approfondissement — Vers la Terminale}\end{center}
% BEGIN-VERIFY
% from sympy import binomial
% for n in range(1,8):
%     assert sum(binomial(n,k) for k in range(n+1)) == 2**n
% END-VERIFY
\begin{corrige}{1SPE-VARALEA-EX-028}

\textbf{1.} $\binom{n-1}{k-1} + \binom{n-1}{k} = \frac{(n-1)!}{(k-1)!(n-k)!} + \frac{(n-1)!}{k!(n-1-k)!}$.

$= \frac{(n-1)! \cdot k + (n-1)! \cdot (n-k)}{k!(n-k)!} = \frac{(n-1)! \cdot n}{k!(n-k)!} = \frac{n!}{k!(n-k)!} = \binom{n}{k}$.

\medskip
\textbf{2.} Par la formule du binôme de Newton avec $a=b=1$ :
$\sum_{k=0}^{n} \binom{n}{k} 1^k \cdot 1^{n-k} = (1+1)^n = 2^n$.

\medskip
\textbf{3.} Pour $X \sim \mathcal{B}(n, \frac{1}{2})$ :
$\sum_{k=0}^{n} P(X=k) = \sum_{k=0}^{n} \binom{n}{k} \left(\frac{1}{2}\right)^n = \frac{1}{2^n} \sum_{k=0}^{n} \binom{n}{k} = \frac{2^n}{2^n} = 1$. \checkmark

\end{corrige}
